\documentclass[11pt]{article}
\usepackage[margin=1in]{geometry}
\usepackage{amsmath,amssymb,mathtools}
\usepackage{microtype}

\newcommand{\E}{\mathbb{E}}
\newcommand{\R}{\mathbb{R}}
\newcommand{\T}{\mathbb{T}}
\newcommand{\Sd}{\mathbb{S}^{d-1}}
\newcommand{\abs}[1]{\lvert #1 \rvert}

\begin{document}

\begin{center}
{\Large Prescribing a power spectrum from Fourier-space variance}
\end{center}

Let $u$ be a statistically isotropic velocity field on $\R^d$ (or on $\T^d$), with Fourier transform $\widehat u(\xi)$. Define the one-dimensional energy spectrum by
\[
E(k)
:= \frac{d}{dk}\int_{\{\abs{\xi}\le k\}} \E\abs{\widehat u(\xi)}^2\,d\xi
= \int_{\{\abs{\xi}=k\}} \E\abs{\widehat u(\xi)}^2\,dS(\xi),
\qquad k>0.
\]
This definition is the same in every dimension. The dimension dependence enters only through the size of the shell $\{\abs{\xi}=k\}$.

Assume the Fourier coefficients are generated by
\[
\widehat u(\xi)=\sigma(\abs{\xi})\,g_\xi,
\qquad \E\abs{g_\xi}^2=1,
\]
so that $\sigma(k)^2=\E\abs{\widehat u(\xi)}^2$ for $\abs{\xi}=k$. Then, on $\R^d$,
\[
E(k)=\abs{\Sd}\,k^{d-1}\sigma(k)^2.
\]
Hence, if one wants
\[
E(k)\sim Ck^{-p},
\]
one should choose
\[
\sigma(k)^2\sim C\,\abs{\Sd}^{-1}k^{-(p+d-1)}.
\]
Equivalently, the standard deviation of each Fourier coefficient on the shell $\abs{\xi}=k$ should scale like
\[
\sigma(k)\sim k^{-(p+d-1)/2}.
\]
On a discrete Fourier grid, the same scaling holds up to constants, since the number of modes in a shell of radius $k$ is $\asymp k^{d-1}$.

\paragraph{Kolmogorov example.}
If $E(k)\propto k^{-5/3}$, then $p=5/3$ and therefore
\[
\sigma(k)^2\propto k^{-(d+2/3)},
\qquad
\sigma(k)\propto k^{-(d/2+1/3)}.
\]
In particular,
\[
\sigma(k)^2\propto k^{-8/3}\quad(d=2),
\qquad
\sigma(k)^2\propto k^{-11/3}\quad(d=3).
\]

\paragraph{Relation with H\"older regularity.}
Heuristically (and exactly in the standard Gaussian scaling models), if the velocity field has H\"older exponent $\alpha$, then
\[
\E\abs{\widehat u(\xi)}^2\sim \abs{\xi}^{-(2\alpha+d)}.
\]
After integrating over the shell $\abs{\xi}=k$, this gives
\[
E(k)\sim k^{d-1}k^{-(2\alpha+d)}=k^{-(2\alpha+1)}.
\]
Thus, if
\[
E(k)\sim k^{-p},
\]
then the corresponding H\"older exponent is
\[
p=2\alpha+1,
\qquad
\alpha=\frac{p-1}{2}.
\]
In particular, the Kolmogorov law $E(k)\propto k^{-5/3}$ corresponds to $\alpha=1/3$.

\medskip
\noindent
All proportionality constants depend on Fourier normalization, shell thickness, and any additional constraints such as reality or divergence-free projection, but the exponents above are unchanged.

\end{document}
